\begin{table}[!htbp]
\centering
\caption{Main model, 9th grade, Portuguese Language: ATT by relative period}
\label{tab:ap_event_base_9th_lp}
\vspace{-0.15cm}
\scriptsize
\setlength{\tabcolsep}{4pt}
\renewcommand{\arraystretch}{0.92}
\resizebox{\textwidth}{!}{%
\begin{tabular}{lcccc}
\toprule
Relative period & TWFE & CS21 & DR-DiD & S\&X approx. \\
\midrule
-4 & 0.084 (0.036) & 0.059 (0.043) & 0.019 (0.060) & 0.047 (0.050) \\
-3 & 0.068 (0.032) & 0.057 (0.033) & 0.016 (0.043) & 0.080 (0.037) \\
-2 & 0.022 (0.027) & 0.022 (0.027) & 0.003 (0.030) & 0.044 (0.031) \\
-1 & 0.000 & 0.000 & 0.000 & 0.000 \\
0 & 0.290 (0.035) & 0.281 (0.037) & 0.303 (0.039) & 0.292 (0.045) \\
1 & 0.396 (0.036) & 0.407 (0.038) & 0.425 (0.038) & 0.416 (0.050) \\
2 & 0.489 (0.037) & 0.507 (0.039) & 0.528 (0.040) & 0.509 (0.056) \\
3 & 0.539 (0.041) & 0.580 (0.044) & 0.598 (0.046) & 0.597 (0.067) \\
4 & 0.661 (0.046) & 0.683 (0.053) & 0.698 (0.051) & 0.695 (0.087) \\
\bottomrule
\end{tabular}
}
\begin{flushleft}
\footnotesize Note: each cell reports the ATT and standard error in parentheses, in standard deviations of annual proficiency. Period \(k=-1\) is normalized as the baseline and therefore appears without a standard error.
\end{flushleft}
\end{table}



